\documentclass[lang=cn,a4paper,11pt,openany,scheme=chinese,twoside]{elegantbook}

\let\arrowvert\relax
%\usepackage{unicode-math}

\let\arrowvert\relax

\usepackage{ctex}
\usepackage{xeCJK}
\usepackage{amsmath} % 建议
\usepackage{amssymb} % 可选，提供更多符号

\usepackage{ctex}
\let\openbox\relax

\usepackage{tikz}
\usetikzlibrary{decorations.pathreplacing,arrows.meta}

%调整页边距
\usepackage{geometry}
\geometry{top=2cm, bottom=2.25cm, outer=2cm, inner=2.5cm}


\usepackage{setspace}
\setstretch{1.3}

%设定行间距
\usepackage{calc}
\setlength{\lineskip}{\baselineskip-\ccwd}
\setlength{\lineskiplimit}{2.5pt}



%浮动图片
\usepackage{wrapfig}
%图标表格注记
\usepackage{caption}
%圆圈数字
\usepackage{pifont}
%表格换行
\usepackage{makecell}
%表格颜色
\usepackage[table]{xcolor}
\colorlet{rowcolcolor}{structurecolor!25!white}

%行间距加大
\newcommand{\wideline}{\setlength{\lineskip}{\baselineskip-\ccwd}\setlength{\lineskiplimit}{2.5pt}}
%蓝色加粗
\definecolor{titleblue}{HTML}{3498DB}
\newcommand{\bluebf}[1]{\textcolor{structurecolor}{\textbf{#1}}}


%向量
\usepackage[f]{esvect}

\renewcommand{\vector}[1]{\vv{#1}}
%平面
\newcommand{\plane}{\text{平面}\ }
%平行且等于
\newcommand{\paraleq}{%
    \mathrel{\text{%
        \tikz[baseline]
        \draw (.1em,0ex) -- (.9em,0ex)
        (.1em,-.425ex) -- (.9em,-.425ex)
        (.350em,.1ex) -- (.650em,1.5ex)
        (.550em,.1ex) -- (.850em,1.5ex);}}}

\newcommand{\mb}{\mathbb}
\newcommand{\z}{\text}
\renewcommand{\a}{\alpha}
\renewcommand{\b}{\beta}
\renewcommand{\d}{\delta}
\newcommand{\D}{\Delta}
\renewcommand{\l}{\lambda}
\newcommand{\m}{\mu}
\newcommand{\fl}{\fillin}
\newcommand{\pr}{\paren}
\newcommand{\p}{\uppi}
\renewcommand{\th}{\theta}
\newcommand{\lam}{\lambda}
\newcommand{\rttri}{\mathrm{Rt}\triangle}
\newcommand{\npar}{\par\noindent}
\DeclareSymbolFont{ugmL}{OMX}{mdugm}{m}{n}
\DeclareMathAccent{\wideparen}{\mathord}{ugmL}{"F3}
\newcommand{\sigmaalg}{\sigma-\text{代数}}
\newcommand{\setjot}[1][0]{\setlength{\jot}{#1 pt}}


%填空题下划线
\newcommand{\fillin}[1][3.5]{%
  \nolinebreak % 不在此处断行
  \hspace{0.2em}%
  \rule[-0.5ex]{#1em}{0.5pt}% 轻微下沉模拟“下划线”效果
  \hspace{0.05em}
}

%调整环境内图片环绕
\usepackage{mwe}
\newcommand{\wrapfix}[1][-1]{\begin{wrapfigure}{r}{0.001\textwidth}%
    \vspace{#1 em}%
    \includegraphics[width=0.001\textwidth]{example-image}%
    \end{wrapfigure}}


%分式使用行间版
\everymath{\displaystyle}


\usepackage{myenvironment}


%elegantbook部分修改
\usepackage{elegantfix}

%强制右页开始
\makeatletter
\newcommand{\ForceRightPage}{%
  \clearpage
  \ifodd\value{page}\relax
    % 已经是奇数页，什么也不做
  \else
    \hbox{}\thispagestyle{empty}\newpage
  \fi
}
\makeatother

\raggedbottom


\begin{document}


\title{动力系统笔记}
\author{庞逸天}
\date{}
\maketitle
\setcounter{page}{1}

\frontmatter
\pagenumbering{roman}     % 目录等用罗马数字
\clearpage
\tableofcontents
\clearpage

\ForceRightPage
\mainmatter

\chapter{预备知识}

\section{概率测度空间}

\begin{definition}{$\sigma-$代数}
    设$X$是一个集合，$\mathcal{B}$是$X$的一些子集构成的集合类，若$\mathcal{B}$满足以下3个条件：
    \begin{enumerate}
        \item $X\in\mathcal{B}$
        \item $A\in\mathcal{B}\Rightarrow X\setminus A\in\mathcal{B}$
        \item $A_1,A_2,A_3,\cdots,A_n,\cdots\in\mathcal{B}\Rightarrow\bigcup_{i=1}^{\infty}{A_i\in\mathcal{B}}$
    \end{enumerate}
    则称$\mathcal{B}$是$X$的一个$\sigma-\text{代数}$
\end{definition}

除了上述定义中满足的三个条件之外，还可以证明$\sigma-\text{代数}$满足以下两个性质.

\begin{proposition}{$\sigma-\text{代数}$的基本性质}
    \begin{enumerate}
        \item $\varnothing\in\mathcal{B}$
        \item 设$A_i\in\mathcal{B}(i=1,2,\cdots,n)$，则$\bigcap_{i=1}^{\infty}{A_i}\in\mathcal{B}$   
    \end{enumerate}
\end{proposition}
\begin{proof}
    \begin{enumerate}
        \item 由条件1可知$X\in\mathcal{B}$，所以$X\setminus X\in\mathcal{B}$，即$\varnothing\in\mathcal{B}$.
        \item 由$A_i\in\mathcal{B}$知$A_i^c=X\setminus A_i \in\mathcal{B}$.所以$\bigcup_{i=1}^\infty{A_i^c}\in\mathcal{B}$
        $\Rightarrow X\setminus\left(\bigcup_{i=1}^\infty{A_i^c}\right)=\bigcap_{i=1}^\infty{A_i}\in\mathcal{B}$.
    \end{enumerate}
\end{proof}

因此，$\sigma-$代数是包含全集合$X$和空集$\varnothing$，且对取余，可数并和可数交运算均封闭的集类.我们也把$(X,\mathcal{B})$称为一个\bluebf{可测空间}，把$(X,\mathcal{B})$中的元素叫做\bluebf{可测集}.下面我们引入可测空间上的测度.

\begin{definition}{概率测度}
    \wideline
    设$X$是一个集合，$\mathcal{B}$是$X$的一个$\sigma-\text{代数}$，函数：$m:\mathcal{B}\rightarrow[0,1]$.如果$m$满足以下2个条件：

    \begin{enumerate}
        \item $m(\varnothing)=0,m(X)=1$
        \item $\forall$两两不交集合$\{A_i\}_{i\in\mathbb{N}}(A_i\in\mathcal{B})$，$m\left(\bigcup_{i=1}^{\infty}{A_i}\right)=\sum_{i=1}^{\infty}{m(A_i)}$
    \end{enumerate}
    则称$m$为$(X,\mathcal{B})$上的一个概率测度，$(X,\mathcal{B},m)$是一个概率测度空间。
\end{definition}
特别地，如果在上面的定义中去掉$m(X)=1$的限制，则称$m$为\bluebf{有限测度}，$(X,\mathcal{B},m)$为\bluebf{有限测度空间}.

\begin{instance}
    (1)取$X=[0,1]$，$\mathcal{B}=\{X\text{中的可测集}\}$，由可测集的性质可知$\mathcal{B}$是$X$的一个$\sigma$-代数.\par

    再令$m=m_{\mathrm{Leb}}:\mathcal{B}\rightarrow [0,1]$，其中$m_{\mathrm{Leb}}$为$\mathbb{R}$上的Lebesgue测度.\par
    由$m_{\mathrm{Leb}}(\varnothing)=0,m_{\mathrm{Leb}}([0,1])=1$以及Lebesgue测度的可列可加性可知$m_{\mathrm{Leb}}$是$(X,\mathcal{B})$的一个概率测度，$(X,\mathcal{B},m_{\mathrm{Leb}})$是一个概率测度空间.\par
    如果把$X$拓展到实数域$\mathbb{R}$，则$m_{\mathrm{Leb}}$就是一个有限测度，$(X,\mathcal{B},m_{\mathrm{Leb}})$是一个有限测度空间.\par
    (2)取$X=\{1,2,\cdots,n\}$，$\mathcal{B}=\{X\text{中的所有子集}\}$，同样可证$\mathcal{B}$是$X$的一个$\sigma$-代数.\par
    定义$m:{\setlength{\jot}{-0.5em}\begin{aligned}
        \mathcal{B}&\to [0,1]\\
        A&\mapsto \frac{|A|}{|X|}
    \end{aligned}}$，则有$m(\varnothing)=\frac{0}{n}=0$，$m(X)=\frac{n}{n}=1$.且若$A_1,A_2,\cdots,A_m$为$X$中两两不交子集，则$m\left(\bigcup_{i=1}^m{A_i}\right)=\frac{\left|\displaystyle\bigcup_{i=1}^m{A_i}\right|}{|X|}=\frac{\displaystyle\sum_{i=1}^m{|A_i|}}{|X|}=\sum_{i=1}^m{\frac{|A_i|}{|X|}}=\sum_{i=1}^m{m(A_i)}$.\par
    因此$m$是$(X,\mathcal{B})$上的概率测度，$(X,\mathcal{B},m)$是一个概率测度空间.\par
    如果我们把$X$看作一个样本空间，$m$实际上描述的就是古典概型中概率的定义.
\end{instance}

但很多时候，$\sigma-\text{代数}$上的测度不容易计算，因此我们需要引入\bluebf{半代数}与\bluebf{代数}的概念.
\begin{definition}{半代数}
    设$X$是一个集合，$\mathcal{S}$是由$X$的一些子集构成的集合类，如果$\mathcal{S}$满足以下3个条件：
    \begin{enumerate}
        \item $\varnothing\in\mathcal{S}$
        \item $A,B\in\mathcal{S}\Rightarrow A\cap B\in\mathcal{S}$
        \item $A\in\mathcal{S}\Rightarrow \exists\text{有限多个两两不交集合}E_i\in\mathcal{S}(i=1,2,\cdots,n)$，使得$X\setminus A=\bigcup_{i=1}^n{E_i}$
    \end{enumerate}
    称$\mathcal{S}$是$X$的一个半代数.
\end{definition}

\begin{instance}
    (1)取$X=[0,1]$，$\mathcal{S}=\{X\text{中所有区间}\}\cup\{\varnothing\}$.
    则$\mathcal{S}$是$X$的一个半代数.\par
    \begin{tikzpicture}[x=8cm,y=1cm,>=Latex,
  tick/.style={thick},
  braceA/.style={decorate,decoration={brace,amplitude=5pt},ultra thick,red!75},
  braceB/.style={decorate,decoration={brace,mirror,amplitude=4pt},ultra thick,blue!70},
  braceE/.style={decorate,decoration={brace,amplitude=5pt},ultra thick,orange!85!brown}
]

% ---------- 左图 ----------
% 数轴
\draw[tick] (-0.05,0)--(0.75,0) [->];
% 刻度与标签
\draw[tick] (0,0.08)--(0,-0.08) node[below=3pt] {$0$};
\draw[tick] (0.7,0.08)--(0.7,-0.08) node[below=3pt] {$1$};

% 区间端点（可根据需要调整）
\def\xA{0.15}  % A 的左端
\def\yA{0.00}
\def\xAmid{0.30} % A 与 B 的交叠开始
\def\xBend{0.50} % B 的右端

% A（红色，大括号）
\draw[braceA] (\xA,0.15) -- (0.4,0.15)
  node[midway,above=5pt,red!80] {$A$};
% B（蓝色，大括号）
\draw[braceB,blue] (0.3,-0.15) -- (0.5,-0.15)
  node[midway,below=5pt,blue] {$B$};

% 视觉上的端点竖线（与手绘一致，可选）
\draw[thick,red!75] (\xA,0.08)--(\xA,-0.08);
\draw[thick,red!75] (0.4,0.08)--(0.4,-0.08);
\draw[thick,blue!70] (0.3,0.08)--(0.3,-0.08);
\draw[thick,blue!70] (0.5,0.08)--(0.5,-0.08);

% 数学公式
\node at (0.35,-1.25) {$A,B\in\mathcal{S}\Rightarrow A\cap B\in\mathcal{S}$}

% ---------- 右图 ----------
\begin{scope}[xshift=6.5cm]
  % 数轴
  \draw[tick] (0.15,0)--(0.95,0) [->];
  % 刻度
  \draw[tick] (0.2,0.08)--(0.2,-0.08) node[below=3pt] {$0$};
  \draw[tick] (0.9,0.08)--(0.9,-0.08) node[below=3pt] {$1$};

  % 端点
  \def\xEoneL{0.2}  % E1 左
  \def\xEoneR{0.4}  % E1 右（也是 A 左）
  \def\xAtwoR{0.7}  % A 右（也是 E2 左）
  \def\xEtwoR{0.9}  % E2 右

  % E1 橙色
  \draw[braceE,orange!85!brown] (\xEoneL,0.15)--(\xEoneR,0.15)
     node[midway,above=6pt,orange!85!brown] {$E_1$};

  % A 红色
  \draw[braceA,red!75] (\xEoneR,0.15)--(\xAtwoR,0.15)
     node[midway,above=6pt,red!75] {$A$};

  % E2 橙色
  \draw[braceE,orange!85!brown] (\xAtwoR,0.15)--(\xEtwoR,0.15)
     node[midway,above=6pt,orange!85!brown] {$E_2$};

  % 端点竖线（可选）
  \foreach \p/\c in {\xEoneL/orange!85!brown,\xEoneR/red!75,\xAtwoR/orange!85!brown,\xEtwoR/orange!85!brown}{
    \draw[thick,\c] (\p,0.08)--(\p,-0.08);
  }

  \node at (0.55,-1.25) {$A\in\mathcal{S}\Rightarrow X\setminus A=E_1\cup E_2$}
\end{scope}
\end{tikzpicture}\par
(2)取$X=[0,1]^n$，$\mathcal{S}=\{X\text{中的所有矩体}\}\cup\{\varnothing\}$，则$\mathcal{S}$是$X$的一个半代数.
\end{instance}



\begin{definition}{代数}
    设$X$是一个集合，$\mathcal{A}$是$X$的一些子集构成的集合类，如果$\mathcal{A}$满足以下3个条件：
    \begin{enumerate}
        \item $\varnothing\in\mathcal{A}$
        \item $A,B\in\mathcal{A}\Rightarrow A\cap B\in\mathcal{A}$
        \item $A\in\mathcal{A}\Rightarrow A^c=X\setminus A\in\mathcal{A}$
    \end{enumerate}
    称$\mathcal{A}$是$X$的一个代数.
\end{definition}
结合第一个和第二个条件，还可以得到代数的以下性质：
\begin{proposition}{代数的性质}
    若$\mathcal{A}$是$X$的一个代数，则若$A_1,A_2,\cdots,A_n\in\mathcal{A}$，$\bigcup_{i=1}^n{A_i},\bigcap_{i=1}^n{A_i}\in\mathcal{A}$.
\end{proposition}
这说明代数是对取余、有限并和有限交封闭的集类，它和$\sigma-$代数的区别在于它不要求对可数并和可数交运算封闭.\par
从半代数的例子中可以看出，$[0,1]^n$内的矩体和空集构成半代数，相较于所有的可测集构成的$\sigma-\text{代数}$要简单许多，因此研究清楚半代数并将这些结论延申到$\sigma-$代数上也会比直接在$\sigma-$代数上做文章方便，为此，需要介绍\bluebf{最小代数}和\bluebf{延拓定理}.\par
\begin{definition}{最小代数}
    设$X$是一个集合，$\mathcal{S}$是$X$的一个半代数，定义\[\mathcal{A}(\mathcal{S})=\bigcap_{\substack{\mathcal{A}\supseteq\mathcal{S}\\\mathcal{A}\,\text{代数}}}{\mathcal{A}}\]
    则$\mathcal{A}(\mathcal{S})$是一个代数，称为包含$\mathcal{S}$的最小代数.
\end{definition}
下面我们证明$\mathcal{A}(\mathcal{S})$是一个代数.
\begin{proof}
    (1)$\varnothing\in\mathcal{A}(\mathcal{S})$\par
    (2)$A,B\in\mathcal{A}(\mathcal{S})\Rightarrow A,B\in\mathcal{A}(\mathcal{A}\subseteq\mathcal{S}) \Rightarrow A\cap B\in\mathcal{A}\Rightarrow A\cap B\in\mathcal{A}(\mathcal{S})$\par
    (3)$A\in\mathcal{A}(\mathcal{S})\Rightarrow A\in\mathcal{A}(\mathcal{A}\supseteq\mathcal{S}) \Rightarrow X\setminus A\in\mathcal{A}(\mathcal{A}\subseteq\mathcal{S})\Rightarrow X\setminus A\in\mathcal{A}(\mathcal{S})$
\end{proof}
\begin{remark}
    可以发现，在上述证明过程中并未使用到$\mathcal{S}$为半代数的性质，因此上述定义实际上并不依赖于$\mathcal{S}$为半代数.
\end{remark}
类似地，可以定义最小$\sigma-\text{代数}$的概念.
\begin{definition}{最小$\sigmaalg$}
    设$X$是一个集合，$\mathcal{S}$是$X$的一个半代数，定义\[\mathcal{B}(\mathcal{A})= \bigcap_{\substack{\mathcal{B}\supseteq\mathcal{A}\\\mathcal{B}\,\sigmaalg}}{\mathcal{B}}\]
    则$\mathcal{B}(\mathcal{A})$是一个$\sigmaalg$，称为包含$\mathcal{A}$的最小$\sigmaalg$.
\end{definition}
同样可以证明$\mathcal{B}(\mathcal{A})$为一个$\sigmaalg$.
\begin{proof}
    (1)$X\in\mathcal{B}(\mathcal{A})$\par
    (2)$A\in\mathcal{B}(\mathcal{A})\Rightarrow A\in\mathcal{B}(\mathcal{B}\supseteq\mathcal{A})\Rightarrow X\setminus A\in\mathcal{B}(\mathcal{B}\supseteq\mathcal{A})\Rightarrow X\setminus A\in\mathcal{B}(\mathcal{A})$\par
    (3)$A_1,A_2,\cdots,A_n,\cdots\in\mathcal{B}(\mathcal{A})\Rightarrow A_1,A_2,\cdots,A_n,\cdots\in\mathcal{B}(\mathcal{B}\supseteq\mathcal{A})\Rightarrow \bigcup_{i=1}^\infty{A_i}\in\mathcal{B}(\mathcal{B}\supseteq\mathcal{A}) \Rightarrow\bigcup_{i=1}^\infty{A_i}\in\mathcal{B}(\mathcal{A})$
\end{proof}
\begin{remark}
    同样地，上述定义也不依赖于$\mathcal{A}$是一个代数.
\end{remark}
除了上述的基本定义之外，我们发现可以用半代数的结构描述其最小代数.
\begin{theorem}
    设$X$是一个集合，$\mathcal{S}$是$X$的一个半代数，则$\mathcal{S}$生成的最小代数$\mathcal{A}(\mathcal{S})$的每个子集均能表示成$\mathcal{S}$中有限个互不相交元素的并，即：
    \[\mathcal{A}(\mathcal{S})=\left\{E=\bigcup_{i=1}^n{E_i}\bigg| E_i\in\mathcal{S},\,E_i\text{两两不交},\,n\in\mathbb{R} \right\}\]
\end{theorem}
\begin{proof}
    记$\mathcal{F}=\mathcal{A}(\mathcal{S})=\left\{E=\bigcup_{i=1}^n{E_i}\bigg| E_i\in\mathcal{S},\,E_i\text{两两不交},\,n\in\mathbb{R} \right\}$.\par
    \begin{enumerate}[label=(\arabic*)]
        \item $\mathcal{F}\supseteq\mathcal{S}$.\par
        令$n=1$即可.
        \item $\mathcal{F}$是$X$的一个代数
        \begin{enumerate}[label=(\roman*)]
            \item $\varnothing\in\mathcal{F}$
            \item 若$A,B\in\mathcal{F}$，则$A=E_1\cup\cdots\cup E_m(E_i\in\mathcal{S})$，$B=F_1\cup\cdots\cup F_n (F_i\in\mathcal{S})$.\par
            所以
            {\setlength{\jot}{-2pt}\begin{align*}
                A\cap B&=(E_1\cup\cdots\cup E_m)\cap(F_1\cup\cdots\cup F_n)\\
                &=(E_1\cap F_1)\cup(E_2\cap F_1)\cup\cdots\cup (E_m\cap F_1)\cup\cdots\\
                &=\bigcup_{\substack{1\leqslant i\leqslant m\\ 1\leqslant j\leqslant n}}(E_i\cap F_j)
            \end{align*}}
            其中$E_i\cap F_j\in\mathcal{S}$且它们两两不交.
            \item 若$A\in\mathcal{F}$，$A=E_1\cap E_2\cap\cdots\cap E_m(E_i\in\mathcal{S},\,\text{两两不交})$.\par
            则$X\setminus A = X\setminus (E_1\cup E_2\cup\cdots\cup E_m)=E_1^c\cup E_2^c\cup\cdots\cup E_m^c$.\par
            由半代数的性质可知，$E_i^c=\bigcup_{j=1}^{n_i}{F_{ij}}$，所以
            {\setlength{\jot}{-2pt}\setlength\belowdisplayskip{-10pt}\begin{align*}
                X\setminus A&=E_1^c\cup E_2^c\cup\cdots\cup E_m^c\\
                &=\bigcup_{j=1}^{n_1}{F_{1j}}\cap\bigcup_{j=1}^{n_2}{F_{2j}}\cap\cdots\cap\bigcup_{j=1}^{n_m}{F_{mj}} \\
                &=(F_{11}\cap F_{21}\cap F_{31}\cap\cdots\cap F_{m1})\cup\cdots \\
            \end{align*}}
            上式中的每一项都在$\mathcal{S}$中且两两不交.
        \end{enumerate}
    \end{enumerate}\par
    结合以上两点可知$\mathcal{F}\supseteq\mathcal{A}(\mathcal{S})$.
    \begin{enumerate}[label=(\arabic*)]
        \setcounter{enumi}{2}
        \item $F\subseteq\mathcal{A}(\mathcal{S})$\par
        $\forall \bigcup_{i=1}^n{E_i}\subset\mathcal{F},\, E_i\in\mathcal{S},\, E_i\text{两两不交},\,n\in\mathbb{N}$，$E_i\subseteq\mathcal{S}\subseteq\mathcal{A}(\mathcal{S})\Rightarrow \bigcup_{i=1}^n{E_i}\in\mathcal{A}(\mathcal{S})$.\par
        所以$\mathcal{A}(\mathcal{S})\supseteq\mathcal{F}$.
    \end{enumerate}\par
    综上可得$\mathcal{F}=\mathcal{A}(\mathcal{S})$.
\end{proof}
在建立起由半代数构造其最小代数的方法后，就可以把半代数上的某些函数延拓到其最小代数上，\bluebf{有限可加函数}是其中的一类.
\begin{definition}{有限可加(可列可加)函数}
    \wideline
    设$X$是一个集合，$\mathcal{S}$是一个半代数，$\tau:\mathcal{S}\to[0,1]$，如果
    \begin{enumerate}
        \item $\tau(\varnothing)=0$
        \item $\forall\bigcup_{i=1}^n{E_i}\in\mathcal{S},\, E_i\in\mathcal{S}(i=1,2,\cdots,n),\, E_i两两不交$，有$\tau\left(\bigcup_{i=1}^n{E_i}\right)=\sum_{i=1}^n{\tau(E_i)}$
    \end{enumerate}
    则称$\tau$是有限可加的.\par
    如果$n=\infty$，称$\tau$是可列可加的.
\end{definition}
显然，可列可加函数包含有限可加函数.\par
在代数$\mathcal{A}$和$\sigmaalg\mathcal{B}$上，也可以类似地定义有限可加函数和可列可加函数.例如概率测度空间$(X,\mathcal{B},m)$中的测度$m$是定义在$\sigmaalg\mathcal{B}$上的可列可加函数.\par
下面介绍半代数到其最小代数的延拓定理.\par
\begin{theorem}
    设$X$是一个集合，$\mathcal{S}$是一个半代数，$\tau:\mathcal{S}\to[0,1]$是一个有限可加(可列可加)函数，那么$\tau$可以唯一地延拓到$\mathcal{A}(\mathcal{S})$上的一个有限可加(可列可加)函数$\tau_1$.(即$\forall E\in\mathcal{S},\tau(E)=\tau_1(E)$)
\end{theorem}
\begin{proof}
    设$\tau:\mathcal{S}\to[0,1]$是一个有限可加函数，定义
    {\setlength{\jot}{-5pt}\begin{align*}
        \tau_1:\mathcal{A}(\mathcal{S})&\to[0,1]\\
        \bigcup_{i=1}^n{E_i}&\to\sum_{i=1}^n{\tau(E_i)}\,\,(E_i\in\mathcal{S},\,\text{两两不交})
    \end{align*}}
    \begin{itemize}
        \item 函数$\tau_1$是良定的.\par
        若$\bigcup_{i=1}^n{E_i}=\bigcup_{j=1}^m{F_j}\in\mathcal{A}(\mathcal{S})(E_i,F_j\in\mathcal{S})$，则
        {\setjot[-5]\begin{align*}
            \tau(E_i)&=\tau\left(E_i\cap\left(\bigcup_{j=1}^m{F_j}\right)\right)=\tau((E_i\cap F_1)\cup(E_i\cap F_2) \cup \cdots\cup(E_i\cap F_m) \\
            &=\tau(E_i\cap F_1)+\tau(E_i\cap F_2)+\cdots+\tau(E_i\cap F_m)=\sum_{j=1}^m\tau(E_i\cap F_j) \\
            \Rightarrow\sum_{i=1}^n\tau&(E_i)=\sum_{i=1}^n\sum_{j=1}^m{\tau(E_i\cap F_j)}
        \end{align*}}
        类似地，$\sum_{j=1}^m{\tau(F_j)}=\sum_{j=1}^m\sum_{i=1}^n{\tau(F_j\cap E_i)}$.\par
        所以$\sum_{i=1}^n{\tau(E_i)}=\sum_{j=1}^m{\tau(F_j)}$，$\tau_1$良定.
        \item 函数$\tau_1$有限可加.\par
        \begin{enumerate}[(1)]
            \item $\tau_1(\varnothing)=\tau(\varnothing)=0$
            \item 若$A_i\in\mathcal{A}(\mathcal{S})(i=1,2,\cdots,n)$两两不交.\par
            由代数的定义可知$A_i=\bigcup_{j=1}^{k_i}{E_{ij}},\,E_ij\in\mathcal{S},\,E_{ij}\text{两两不交},\, 1\leqslant j\leqslant k_i$.\par
            所以
            \vspace{-1em}
            {\setjot[-5]\begin{align*}
                \sum_{i=1}^n{\tau_1(A_i)}&=\tau_1(A_1)+\tau_1(A_2)+\cdots+\tau_1(A_n)\\
                &=\sum_{j=1}^{k_1}{\tau(E_{1j})}+\sum_{j=1}^{k_2}{\tau(E_{2j)}}+\cdots+\sum_{j=1}^{k_n}{\tau(E_{nj)}}\\
                &=\sum_{i=1}^n\sum_{j=1}^{k_i}{\tau(E_{ij})}=\tau\left(\bigcup_{i=1}^n\bigcup_{j=1}^{k_i}{E_{ij}}\right)\\
                &=\tau_1\left(\bigcup_{i=1}^n{A_i}\right)
            \end{align*}}
        \end{enumerate}
        \item 函数唯一\par
        如果$\tilde{\tau}$也满足条件，则$\tilde{\tau}=\tau_1$
    \end{itemize}
\end{proof}
类似地有代数$\mathcal{A}$到其最小$\sigmaalg\mathcal{B}(\mathcal{A})$的延拓定理.
\begin{theorem}
    设$X$是一个集合，$\mathcal{A}$是一个代数，$\tau:\mathcal{A}\to[0,1]$是$\mathcal{A}$上的一个可列可加函数，那么$\tau$可以唯一地延拓到$\mathcal{B}(\mathcal{A})$上的可列可加函数$m:\mathcal{B}(\mathcal{A})\to[0,1]$.(即$\forall A\in\mathcal{A},m(A)=\tau(A)$)
\end{theorem}
\begin{instance}
    设$X=[0,1]$，$\mathcal{S}=\{X\text{的所有子区间}\}$为一个半代数，$\mathcal{B}=\{X\text{中的可测集}\}$是一个$\sigmaalg$.\par
    在$\mathcal{S}$上定义函数$\tau:{\setjot[-7]\begin{aligned}
        \mathcal{S}&\to[0,1]\\
        I&\to|I|
    \end{aligned}}$，则$\tau$可延拓到$\mathcal{B}\left(\mathcal{A}(\mathcal{S})\right)$上的概率测度，其中$\mathcal{B}\left(\mathcal{A}(\mathcal{S})\right)$和$\mathcal{B}$只相差一个零侧集.
\end{instance}


\end{document}